\section{Synthesis and Research Agenda}\label{sec:agenda}
\textbf{The field map turns a privacy score into a protection decision by binding its observations to executable work.}
The strongest cross-study consequence is intervention matching: WRA's visibility utility cannot establish executed-decoy utility, and adapted OHTTP inference with traffic cost does not establish original-task preservation. The technical case adds the dispatcher: controller cover can omit model waits and become removable, history matching hides the shortcut but not the signal, and planner cover is bounded by a pair-ambiguity floor that only more model work can lower. These connections extend the evaluation object; the underlying adaptive-attack and state-correctness principles are established.

\subsection{Source dependence}\label{sec:sensitivity}
The evidence map follows 33 statements in seven comparisons across 20 source-removal scenarios (Supplement S4.4). Strictly retaining inspected published texts leaves WRA's tool-destination recovery and decoy/visibility distinction, plus CHeaT's defined detection/prevention endpoints. The model-view observations, live OHTTP refitting and each named Agent transfer axis depend on preprints. The local cover case adds implementation-specific evidence, not independent confirmation of those source results. This early SoK is useful for fixing interpretation and evaluation conditions; the small direct corpus limits field-wide empirical conclusions.

\subsection{Tests that would change the decision}
Table~\ref{tab:open-questions} gives discriminating tests rather than a list of desirable benchmarks. Four follow directly from the live case. (i) A shadow sweep over $k\in\{1,2,3\}$ and over iid versus distinct label sampling should trace the floor of Table~\ref{tab:shadow-floor} against token cost, with an exact real-chain test on the distinct-label fixtures. (ii) The recipient view: each MCP service already receives the decoy arguments, so a per-service attacker can be evaluated on the existing journals. (iii) Network realism: delay jitter and loss injected at the services, a cross-host deployment and unseen request phrasings would show how much of the timing signal survives outside loopback. (iv) Equal-budget comparisons of controller and planner cover under production state contracts. Business invariants should admit legitimate concurrent updates. Future evidence can support a policy by meeting declared inference and utility thresholds, or reject it by exposing a budget or state violation.
\begin{table*}[t]\centering\small
\caption{Testable open questions derived from the synthesis. Q1 specifies transfer tests. Q2 separates the forwarding mediators; Q3 extends the controlled cover study to application constraints; Q4 extends live detection to new workloads and deployments. Each gives a result that could support or revise the proposed explanation.}\label{tab:open-questions}
\setlength{\tabcolsep}{4pt}
\begin{tabularx}{\textwidth}{@{}P{.22\textwidth}P{.43\textwidth}Y@{}}\toprule
Question and unresolved mechanism & Controlled comparison & Discriminating result \\\midrule
\textbf{Q1. Which signals survive execution changes?} Do attacks identify recurring task structure or features tied to its implementation? & Keep task labels and observation rights fixed. Hold out compositions, coordination and client/scaffold versions separately, then jointly. Compare destination, volume/timing, sequence and recipient-visible feature families under declared frozen and adapted training. & Identify a feature family whose inference survives specified changes, or a failure tied to one changed condition. Report per-axis and joint transfer with attempted-run coverage; a pooled average cannot establish invariance. \\\addlinespace[4pt]
\textbf{Q2. How does feedback reach model traffic?} Which part is content, length or later action structure? & Fix initial instructions; project to category-free task fields. Pair length-matched neutral padding with validity-preserving replacement of supplementary content; hold encoding and follow-up policy fixed where possible. Vary record count separately. Compare prefix, first-feedback and complete model views under tool-only and joint protection. & Additional held-out inference after feedback, with controls that exclude the competing path, supports that attribution. Report an effect interval against a prespecified meaningful-effect threshold. A nonsignificant difference alone does not rule out leakage. \\\addlinespace[4pt]
\textbf{Q3. Can cover preserve task semantics?} When should it stop at transport or execute in isolation? & Compare planner-sampled decoy schedules (destination sequence and timing) with marginally matched dispatch under production tool contracts and legitimate concurrent updates. Compare no cover, transport termination and isolated execution under declared resource budgets. Include timing-informed removal and independently captured record times; measure inference, delay, original-task completion, quota, state and planner-visible feedback. & A privacy--utility operating point satisfying prespecified state and task tolerances supports the design. Resource exhaustion, changed business state or impaired completion identifies the constraint it violates. \\\addlinespace[4pt]
\textbf{Q4. Which inference survives task discovery?} Does cross-link association recover information hidden by waits and multiplexing? & Train on one concurrency/load and test new compositions, loads and gateways. Compare tool-only and associated tool/model views with observable flow identifiers; keep ground-truth task and stream mappings outside attacker inputs. Include real unrelated traffic and constant-output detector references. & Measure interval recall, false alarms per unit time and correctly localized inference together. Gains must survive changed workloads and loads at a stated false-alarm budget. \\
\bottomrule\end{tabularx}\end{table*}


\section{Validity, Ethics and Availability}\label{sec:validity}
\textbf{The conclusions apply to the recorded sources and declared observation conditions.}
Coverage combines bounded retrieval with citation tracing; failed indices limit recall. Most direct sources are preprints. The independent coder's decision-critical subset agrees on the four-setting privacy denominator and its adaptation/utility fields; other extracted fields do not inherit that reliability estimate.

The live case runs on one host over loopback, with synthetic tools, two fixed 8B models and fixtures and phrasings reused across splits and studies. Loopback timing is cleaner than a wide-area path, so the timing attacks may be optimistic for a remote observer; network jitter, cross-host deployment and unseen phrasings are untested. Same-service calls serialize, and unknown-task discovery and unrestricted multiplexing are not evaluated. Llama's no-cover goal and state counts (82/120 and 70/120) limit business-equivalence conclusions for that planner. Persistent caches and phase order affect the costs.

\paragraph{Ethics and assistance}
All experimental inputs are synthetic; capture filters restrict collection to the experiment. The adverse-write control runs on isolated fixtures. No customers or production business states are affected. An earlier hosted retest used account-holder authorization; provider research terms remain unverified and its results support no main claim. AI assisted extraction, implementation and writing; measurements are backed by saved observations and executable analysis. Traffic inference can enable profiling. Account credentials, hosted-retest captures and transcripts, and third-party full texts are excluded from distribution.

\paragraph{Availability}
All experimental inputs are synthetic, so the artifact can be complete. It contains the collection and analysis code; fixture generators and seeds; frozen plans, validation selections and their hashes; per-episode record metadata, predictions and call truth; the packet captures with their session keys, which belong only to the synthetic test deployment; the forty-four-setting extraction, source locators, blind re-coding sheets and adjudication record; and scripts that rebuild every table from the records. Model weights are public and are identified by layer hashes. The artifact is available at \artifacturl.

\section{Conclusion}
Executed cover creates both network observations and application work. This SoK maps the existing evidence to those two outputs and provides a joint protocol for residual inference, task completion, business state and resource cost. The live case shows how a controller's dispatch process exposes cover, why history matching does not remove the signal, how a single shadow planner's residual is explained by its pair-ambiguity floor, and why a goal score alone cannot certify the whole execution. Protection claims become actionable when their privacy and business evidence describe the same policy on the same workload.
